\subsection{PR-only execution with a bounded endpoint budget}
\label{sec:pr-engine}
The negative aggregate-cost result for the preceding racer motivated a
disclosed implementation follow-up. The earlier engine crosses Python and
NumPy boundaries during refinement, constructs ROC coordinates even for PR
selection, and can recompute evicted endpoint curves. Saving distinct width
queries therefore need not save actual curve calculations or time.

The follow-up executes the controller and PR curve calculations in compiled
code. A single quantile traversal maintains separate minimum and maximum
deques to compute both p-box extrema, avoiding materialized breakpoint and
coordinate-envelope arrays. It computes the same real-arithmetic bounds;
this is an execution change, not a new p-box theorem. Each observed width's
scalar is retained. Endpoint coordinates are never evicted during adaptive
refinement: before their budget is exhausted, that model switches to an
exhaustive sweep of only its unobserved widths, reusing preparation and
scratch storage. No public exact-evaluator call restarts the computation.

Let $W=L+1$, $G_m$ be model $m$'s number of distinct support-score groups,
and $Q$ be the maximum-support size. The endpoint capacity is
\[
 C_m=\min\!\left(W,\max(2,\lfloor\alpha W\rfloor),
             \left\lfloor\frac{B}{16G_m}\right\rfloor\right),
\]
with $C_m=0$ when $\alpha=0$. If a model's initial merge-block endpoints
already exceed this capacity, it uses exact completion immediately. The
declared initial policy uses $\alpha=1/8$ and $B=8$ MiB per model. Two
Float64 coordinate arrays per retained endpoint account for the factor 16;
$B$ does not bound process RSS or all auxiliary memory.

Each model performs at most $W$ actual width evaluations, including exact
completion, because an observed scalar is never evaluated again. The
artifact separately counts adaptive and completion evaluations. In addition
to prepared inputs, retained coordinates use at most $\sum_m16C_mG_m$
bytes, with $O(MW)$ scalar and controller state. The implementation scans
the model-by-width state when updating global bounds and scans interval
priorities during refinement. Including the initial state scans, its bookkeeping
can require $O(MW(1+\sum_m C_m))$ work, bounded above by $O(M^2W^2)$ when capacity scales
with width. Consequently the endpoint budget is a query/storage guarantee,
not a constant-factor wall-clock overhead guarantee.

An equally specialized exhaustive PR-only sparse kernel and an independent
PR-only dense kernel are required controls. The former uses the same
unobserved-width completion primitive with every width initially unknown;
the latter constructs the sum-and-clip ramps directly. Both execute all
widths in one compiled call and reuse scratch arrays. Thus any advantage of
the new racer must survive comparison with an exact method that also omits
ROC work. Before the policy holdout, two further parameter-free controls
receive the same global radix ranking: one retains direct ramp construction,
and one uses the same linear nearest-interval geometry and weight formula
as the sparse method. Both still traverse every input position and complete
every numerical tie at every width, computing PR alone. These controls
separate generic radix and geometry gains from sparse event elimination;
the nearest variant takes $O((L+1)n)$ time and $O(n+K)$ auxiliary memory.
Float64 guards, separation margins and exact-tie conventions are
unchanged. This follow-up does not replace the earlier negative results.
